% ===================================================================
%  TABLICA Nr 6 — Wnętrze domu, umeblowanie, żywienie, oświetlenie.
%  Meme regime que les precedents.
%
%  LE RENVOI (16) DU BLOC t06-c1-06-1 EST CELUI DU FRANCAIS, NON CELUI
%  DE L'IDO. Le fac-simile ido y ouvre sur « (6) », qui est le savon de
%  l'alinea 2 ; c'est la femme de chambre, « (16) ». C'est l'un des
%  trois ecarts declares de outils/renvoji.py, et les traductions le
%  rendent — le neerlandais et le tcheque ecrivent (16), cette colonne
%  aussi.
%
%  LE VOCABULAIRE DE LA MAISON EST CELUI QUI VIEILLIT LE PLUS VITE, et
%  c'est ici qu'il faut se garder de moderniser les CHOSES. Le
%  « vekigilo » est un reveil a remonter, la « bujio » une bougie, le
%  « petrol-lampo » une lampe a petrole : on ecrit budzik, świeca,
%  lampa naftowa, et non un mot d'aujourd'hui pour un objet de 1926.
%  On modernise la langue, jamais les choses.
% ===================================================================

%%K t06-apar-1 apar
\VUsaut{6.0mm}
\VUtitre{15.0pt}{60}{TABLICA Nr 6}
\VUsaut{1.4mm}
\VUtitre{11.4pt}{0}{Wnętrze domu, umeblowanie, żywienie,}
\VUsaut{0.4mm}
\VUtitre{11.4pt}{0}{oświetlenie.}
\VUsaut{2.2mm}

%%K t06-apar-2 apar
Dom jest ukończony. Wuj Jana zamieszkał w swoim nowym domu, którego rozkład znamy. Teraz
zobaczmy, jak go umeblował. Zwiedzimy najpierw:

%%K t06-tit-1 sub
\VUpk{10.2pt}{40}{Sypialnia.}
\VUsaut{3.0mm}

%%K t06-c1-01-1 p
1. --- Przychodzimy nie w porę, bo pokojówka jest zajęta sprzątaniem \VUgras{pokoju}.

%%K t06-c1-02-1 p
2. --- Na \VUgras{toaletce} \textsuperscript{(1)} są tu i ówdzie rozmaite przedmioty,
którymi się myjemy, czeszemy, krótko mówiąc robimy toaletę. Są to \VUgras{miednica}
\textsuperscript{(2)} i \VUgras{dzbanek na wodę} \textsuperscript{(3)},
\VUgras{szczotka do włosów} \textsuperscript{(4)}, \VUgras{grzebień}
\textsuperscript{(5)}, \VUgras{mydło} \textsuperscript{(6)} i \VUgras{gąbka}
\textsuperscript{(7)}, wreszcie \VUgras{flakonik} \textsuperscript{(8)} na płyn do
zębów.

%%K t06-c1-03-1 p
3. --- Na podłodze, przed toaletką, oto \VUgras{wiadro} \textsuperscript{(9)} do
zbierania brudnej wody.

%%K t06-c1-04-1 p
4. --- Widzimy w \VUgras{przedpokoju} \textsuperscript{(10)} kilka \VUgras{wieszaków}
\textsuperscript{(13)} i dwa \VUgras{parasole} \textsuperscript{(11)} w \VUgras{stojaku
na parasole} \textsuperscript{(12)}.

%%K t06-c1-05-1 p
5. --- Po drugiej stronie drzwi jest \VUgras{szafa z lustrem} \textsuperscript{(14)},
która mieści \VUgras{bieliznę} \textsuperscript{(15)}.

%%K t06-c1-06-1 p
6. --- Skorzystajmy z tego, że \VUgras{pokojówka} \textsuperscript{(16)} ściele
\VUgras{łóżko} \textsuperscript{(17)}, żeby obejrzeć jego rozmaite części.
\VUgras{Rama łóżka} \textsuperscript{(20)} mieści dobry \VUgras{materac sprężynowy}
\textsuperscript{(21)}, na którym Justyna zaraz położy wełniany \VUgras{materac}
\textsuperscript{(22)}. Potem weźmie z krzeseł dwa \VUgras{prześcieradła}
\textsuperscript{(23)} i \VUgras{koc} \textsuperscript{(24)}. Wtedy położy u wezgłowia
\VUgras{poduszkę wałek} \textsuperscript{(25)} i \VUgras{poduszkę}
\textsuperscript{(26)}, które są razem z \VUgras{pierzyną} \textsuperscript{(27)} na
\VUgras{dywaniku przy łóżku} \textsuperscript{(32)}. Używa miotełki z piór, żeby
odkurzyć \VUgras{firanki} \textsuperscript{(19)} i \VUgras{baldachim}
\textsuperscript{(18)} i oto część roboty skończona.

%%K t06-c1-07-1 p
7. --- Potem będzie musiała trochę uporządkować \VUgras{stolik nocny}
\textsuperscript{(28)}, nakręcić \VUgras{budzik} \textsuperscript{(29)}, przygotować
\VUgras{lampkę nocną} \textsuperscript{(30)}, zdjąć \VUgras{świecę}
\textsuperscript{(31)}, żeby wyczyścić świecznik.

%%K t06-tit-2 sub
\VUsaut{5.0mm}
\VUpk{10.2pt}{40}{Koszyk z robótką i przybory kominkowe.}
\VUsaut{3.0mm}

%%K t06-c1-08-1 p
8. --- \VUgras{Koszyk z robótką} \textsuperscript{(34)} obok \VUgras{taboretu}
\textsuperscript{(33)} także bardzo potrzebuje uporządkowania. Będzie musiała złożyć
\VUgras{hafty} \textsuperscript{(35)}, zamknąć w \VUgras{stoliku do roboty}
\textsuperscript{(38)} \VUgras{naparstek}, \VUgras{nici} \textsuperscript{(36)},
\VUgras{igły} \textsuperscript{(37)} i \VUgras{nożyczki} \textsuperscript{(39)} i
zamknąć \VUgras{kluczem} \textsuperscript{(40)} wieko stolika.

%%K t06-c1-09-1 p
9. --- Na \VUgras{stoliku jednonożnym} \textsuperscript{(41)} pośrodku Justyna odnowi
wodę w \VUgras{karafce} \textsuperscript{(42)}. Włoży \VUgras{cukier} do
\VUgras{cukiernicy} \textsuperscript{(43)}, wyczyści \VUgras{szczypce do cukru}
\textsuperscript{(44)} i zabierze \VUgras{szczotkę do ubrania} \textsuperscript{(46)}.

%%K t06-c1-10-1 p
10. --- Prawa strona pokoju będzie od niej wymagała mniej pracy, bo \VUgras{szezlong}
\textsuperscript{(47)} i jego \VUgras{poduszka} \textsuperscript{(48)} są na swoim
właściwym miejscu, a ponieważ nie jest zimno, nic nie zostało przesunięte wśród przyborów
\VUgras{kominka} \textsuperscript{(49)}. \VUgras{Łopatka} \textsuperscript{(50)},
\VUgras{szczypce} \textsuperscript{(51)}, \VUgras{wilki kominkowe}
\textsuperscript{(55)}, \VUgras{zapora na popiół} \textsuperscript{(56)} są czyste;
\VUgras{ekran kominkowy} \textsuperscript{(54)} nie został przesunięty, a
\VUgras{skrzynia na drewno} \textsuperscript{(52)} jest dobrze zaopatrzona w
\VUgras{polana} \textsuperscript{(53)}.

%%K t06-noto-1 noto
\VUnotes{143.2mm}{%
(*) Babo = małe dziecko, ang. baby, fr. bébé, wł. bambino.
\par
}

%%K t06-noto-2 noto
\VUnotes{143.2mm}{%
(*) Lambrequino = fr. lambrequin, hiszp. lambrequines, port. lambrequins, nawet niem. i
ang. lambrequins, więc niem. ang. fr. port. hiszp.%
}

%%K t06-c1-11-1 p
11. --- Będzie jednak musiała odkurzyć \VUgras{zegar wahadłowy} \textsuperscript{(57)},
\VUgras{kandelabry} \textsuperscript{(58)} i \VUgras{pudełka na drobiazgi}
\textsuperscript{(59)}, wreszcie wytrzeć \VUgras{lustro} \textsuperscript{(60)}.

%%K t06-c1-12-1 p
12. --- Co do \VUgras{kołyski} \textsuperscript{(61)} dzieciątka (bobasa (*)), ta jest
już gotowa.

%%K t06-tit-3 sub
\VUsaut{5.0mm}
\VUpk{10.2pt}{40}{Salon.}
\VUsaut{3.0mm}

%%K t06-c2-01-1 p
1. --- Teraz oto salon. To bardzo piękny \VUgras{pokój}. Kominek, który znajduje się w
prawym rogu, jest ozdobiony delikatną \VUgras{statuetką} \textsuperscript{(1)}, dwoma
pięknymi \VUgras{wazonami} \textsuperscript{(3)} i udrapowany aksamitnym
\VUgras{lambrekinem} \textsuperscript{(2)} (*).

%%K t06-c2-02-1 p
2. --- Żeby przeszkodzić iskrom z paleniska zapalić dywan, postawiono przed paleniskiem
miedzianą \VUgras{osłonę przeciw iskrom} \textsuperscript{(4)}.

%%K t06-c2-03-1 p
3. --- Obok elegancki pański \VUgras{sekretarzyk} \textsuperscript{(5)} jest otwarty.
Na nim \VUgras{pani domu} \textsuperscript{(23)} pisze swoje listy.

%%K t06-c2-04-1 p
4. --- Okno jest przybrane \VUgras{firankami} \textsuperscript{(6)} z
\VUgras{podwiązaniami} \textsuperscript{(8)} zaczepionymi o \VUgras{haczyki}
\textsuperscript{(7)} i materiałowymi portierami \VUgras{udrapowanymi u góry}
\textsuperscript{(9)}.

%%K t06-c2-05-1 p
5. --- We wnęce okna \VUgras{osłona doniczki} \textsuperscript{(11)} mieści zieloną
\VUgras{roślinę} \textsuperscript{(10)}, wspaniałą palmę, której liście rozciągają się
prawie aż do \VUgras{lampy naftowej} \textsuperscript{(12)}.

%%K t06-c2-06-1 p
6. --- \VUgras{Abażur} \textsuperscript{(13)} lampy przygotowała Maria, czarująca
\VUgras{panienka} \textsuperscript{(14)}, która siedzi przy fortepianie
\textsuperscript{(15)}. Siedząc bardzo prosto na \VUgras{taborecie}
\textsuperscript{(16)}, uczy się utworu. Jest bardzo zdolna i staranna, jak dowodzi jej
\VUgras{szafka na nuty} \textsuperscript{(17)}, która jest doskonale uporządkowana.

%%K t06-tit-4 sub
\VUsaut{5.0mm}
\VUpk{10.2pt}{40}{Urwis.}
\VUsaut{3.0mm}

%%K t06-c2-07-1 p
7. --- Jej brat, Paweł, szuka \VUgras{książki} \textsuperscript{(19)} w
\VUgras{biblioteczce} \textsuperscript{(18)}. Jest \VUgras{młodzieńcem}
\textsuperscript{(20)}, ruchliwym i nieznośnym. Czepiając się biblioteczki, żeby
dosięgnąć książki, która jest za wysoko, ryzykuje strąceniem \VUgras{popiersia}
\textsuperscript{(21)}, które stoi na górze.

%%K t06-c2-08-1 p
8. --- Korzysta z tego, żeby zrobić to głupstwo, że jego matka jest zajęta rozmową z
przyjaciółką, \VUgras{damą} \textsuperscript{(22)}, która przyjechała spędzić kilka dni w
jej domu. Matka siedzi na \VUgras{kanapie} \textsuperscript{(24)} obok \VUgras{fotela}
\textsuperscript{(25)}, a jej przyjaciółka jest na \VUgras{siedzeniu}
\textsuperscript{(27)}, z którego widać tylko \VUgras{oparcie} \textsuperscript{(26)}.

%%K t06-c2-09-1 p
9. --- Wielka \VUgras{rama} \textsuperscript{(30)} wisząca na ścianie mieści
\VUgras{portret} \textsuperscript{(29)} krewnej. W \VUgras{albumie}
\textsuperscript{(32)} położonym na stole zebrane są \VUgras{fotografie}
\textsuperscript{(33)} pozostałych członków rodziny. Dzieci właśnie go przeglądały, bo
\VUgras{krzesła} \textsuperscript{(31)} są jeszcze skupione wokół stołu.

%%K t06-c2-10-1 p
10. --- \VUgras{Chińska waza} \textsuperscript{(34)} z porcelany, która stoi obok
\VUgras{noża do papieru} \textsuperscript{(35)}, została podarowana Marii na jej
imieniny, a \VUgras{parawan} \textsuperscript{(36)}, który zdobi kąt pokoju, przywiózł z
Chin jej wuj.

%%K t06-c2-11-1 p
11. --- Rozweselony pięknymi \VUgras{obrazami} \textsuperscript{(37)} zawieszonymi na
\VUgras{ścianie} \textsuperscript{(42)}, oświetlony \VUgras{żyrandolem}
\textsuperscript{(38)}, którego \VUgras{lampy elektryczne} \textsuperscript{(39)}
wieczorem świecą radośnie, ten salon wygląda bardzo przyjemnie. Miękki \VUgras{dywan}
\textsuperscript{(40)} rozłożony na parkiecie i tak dogodny \VUgras{dzwonek
elektryczny} \textsuperscript{(41)} dopełniają jego wygody.

%%K t06-tit-5 sub
\VUsaut{3.6mm}
\VUpk{10.2pt}{40}{Jadalnia.}
\VUsaut{3.0mm}

%%K t06-c3-01-1 p
1. --- Zadzwonił dzwon na obiad. Cała rodzina, oprócz Marii, jest zebrana wokół stołu.
Jadalnia jest przyjemnie ogrzana przez wyborny \VUgras{piec} \textsuperscript{(1)}
fajansowy, z cienką \VUgras{rurą} \textsuperscript{(2)}.

%%K t06-c3-02-1 p
2. --- Ten pokój nie ma wielu mebli. Wzdłuż ściany wisi, nad \VUgras{stołkiem}
\textsuperscript{(4)}, ozdobna \VUgras{fontanna} \textsuperscript{(3)}. Całkiem blisko
jest \VUgras{kredens} \textsuperscript{(5)}, na którym kładzie się zimne potrawy, dania
deserowe i rozmaite przybory stołowe, których nie potrzeba od razu.

%%K t06-c3-03-1 p
3. --- Tak widzimy na górnej półce \VUgras{pieczonego kurczaka} \textsuperscript{(7)},
\VUgras{rybę} \textsuperscript{(8)} i \VUgras{czarę} \textsuperscript{(10)} z
\VUgras{owocami} \textsuperscript{(9)}. Niżej oto \VUgras{ser} \textsuperscript{(11)},
\VUgras{chleb} \textsuperscript{(12)}, \VUgras{stojak na flakony} \textsuperscript{(13)},
który mieści wszystko, czego trzeba do przyprawienia sałaty: oliwę, ocet, \VUgras{sól}
\textsuperscript{(14)} i \VUgras{pieprz} \textsuperscript{(15)}. Wreszcie na dolnej
półce jest \VUgras{ciasto} \textsuperscript{(17)} obok \VUgras{musztardniczki}
\textsuperscript{(16)}.

%%K t06-c3-04-1 p
4. --- Zegar jadalny, który nazywa się \VUgras{zegarem ściennym} \textsuperscript{(20)},
ma bardzo osobliwy kształt. Jest wprawiony w drewnianą ramę, która mieści ponadto
\VUgras{barometr} \textsuperscript{(19)} i \VUgras{termometr} \textsuperscript{(18)}.

%%K t06-tit-6 sub
\VUsaut{4.6mm}
\VUpk{10.2pt}{40}{Podawanie obiadu.}
\VUsaut{3.0mm}

%%K t06-c3-05-1 p
5. --- Na wszystkich ścianach pokoju jest położona \VUgras{boazeria}
\textsuperscript{(21)} z woskowanego dębu. Materiałowa \VUgras{portiera}
\textsuperscript{(22)} dosyć szczelnie zamyka drzwi, którymi wchodzi się na werandę. Tam
rodzina pójdzie pić kawę po obiedzie. Już \VUgras{filiżanki} \textsuperscript{(24)} i
\VUgras{spodki} \textsuperscript{(25)} są gotowe na \VUgras{kredensie na naczynia}
\textsuperscript{(23)}.

%%K t06-c3-06-1 p
6. --- Nad stołem jest umocowana do sufitu \VUgras{lampa wisząca}
\textsuperscript{(26)}, której bardzo jasna lampa wieczorem oświetla całą jadalnię.

%%K t06-c3-07-1 p
7. --- Teraz obejrzyjmy sam \VUgras{stół jadalny} \textsuperscript{(27)}. Kiedy rodzina
weszła, nakrycie było gotowe. Na \VUgras{obrusie} \textsuperscript{(28)} był położony,
na miejscu każdego biesiadnika, \VUgras{talerz} \textsuperscript{(33)},
\VUgras{serwetka} \textsuperscript{(29)}, \VUgras{łyżka} \textsuperscript{(34)},
\VUgras{widelec} \textsuperscript{(35)}, \VUgras{nóż} \textsuperscript{(36)} i
\VUgras{szklanka} \textsuperscript{(32)}. Były też \VUgras{butelki}
\textsuperscript{(30)} na wino i \VUgras{karafka} \textsuperscript{(31)} na wodę.

%%K t06-c3-08-1 p
8. --- \VUgras{Służący} \textsuperscript{(39)} przyniósł najpierw \VUgras{wazę}
\textsuperscript{(37)}, w której jeszcze paruje trochę zupy. Teraz podaje pierwsze
\VUgras{danie} \textsuperscript{(38)}, danie mięsne. Potem poda te, które już
nazwaliśmy, wreszcie po \VUgras{deserze} skorzysta z tego, że jego państwo przejdą na
werandę, żeby zabrać przybory stołowe i uporządkować jadalnię.

%%K t06-c3-08-2 p
\textit{Uwaga.} --- \textit{Słowa potoczne odnoszące się do żywności są tu wskazane
bardzo skrótowo. Lista będzie uzupełniona na dwóch tablicach « Hotel » (nr 13) i
« Targ » (nr 15).}

%%K t06-tit-7 sub
\VUsaut{4.6mm}
\VUpk{10.2pt}{40}{Kuchnia.}
\VUsaut{3.0mm}

%%K t06-c4-01-1 p
1. --- W kuchni, na którą zaczynamy patrzeć przez uchylone drzwi, widzimy malowniczy
nieporządek. Przed \VUgras{paleniskiem} \textsuperscript{(1)}, w którym trzeszczy
\VUgras{ogień} \textsuperscript{(3)}, ustawiony jest \VUgras{rożen}
\textsuperscript{(5)}. Piękna \VUgras{pieczeń} \textsuperscript{(7)} jest
\VUgras{nadziana na rożen} \textsuperscript{(6)} i dopieka się.

%%K t06-c4-02-1 p
2. --- \VUgras{Dmuchawka} \textsuperscript{(8)} i \VUgras{szufelka do węgla}
\textsuperscript{(9)}, których dopiero co użyto, zostały na podłodze, tak samo jak
\VUgras{polana} \textsuperscript{(10)}.

%%K t06-c4-03-1 p
3. --- Wierzch kominka jest uporządkowany; \VUgras{latarnia} \textsuperscript{(11)},
\VUgras{pudełko na zapałki} \textsuperscript{(13)}, w którym są \VUgras{zapałki}
\textsuperscript{(12)}, \VUgras{świecznik} \textsuperscript{(14)} i \VUgras{lichtarz}
\textsuperscript{(15)} ze \VUgras{świecą} \textsuperscript{(16)} są na swoim miejscu.
Ale wielkie \VUgras{wiadro na węgiel} \textsuperscript{(18)}, pełne \VUgras{węgla}
\textsuperscript{(17)}, które \VUgras{kucharka} \textsuperscript{(23)} dopiero co
napełniła w piwnicy, zostało pośrodku kuchni.

%%K t06-c4-04-1 p
4. --- Doprawdy, biedna kobieta musi się spieszyć, żeby obiad był gotowy w porę. Teraz
jest przy \VUgras{kuchni} \textsuperscript{(19)} i miesza \VUgras{sos}
\textsuperscript{(24)}, którego nie należy za bardzo przypalić.

%%K t06-c4-05-1 p
5. --- W \VUgras{piekarniku} \textsuperscript{(20)} piecze się ciasto, które
przygotowała na podwieczorek dla dzieci. Nad kuchnią wisi \VUgras{warząchew}
\textsuperscript{(21)} i \VUgras{cedzak} \textsuperscript{(22)}.

%%K t06-tit-8 sub
\VUsaut{4.0mm}
\VUpk{10.2pt}{40}{Naczynia.}
\VUsaut{3.0mm}

%%K t06-c4-06-1 p
6. --- \VUgras{Naczynia kuchenne} \textsuperscript{(25)}, zawieszone w głębi pokoju, są
bardzo czyste. Piękne miedziane \VUgras{rondle} \textsuperscript{(26)}, \VUgras{garnek
na mleko} \textsuperscript{(27)}, \VUgras{sitko} \textsuperscript{(28)} i \VUgras{tarka}
\textsuperscript{(29)} zostały starannie wyczyszczone.

%%K t06-c4-07-1 p
7. --- \VUgras{Solniczka} \textsuperscript{(30)} stoi na kredensie obok \VUgras{dzbanka
na herbatę} \textsuperscript{(31)} i \VUgras{butli na oliwę} \textsuperscript{(32)}.
\VUgras{Szuflada} \textsuperscript{(33)} zapewne mieści sztućce i noże kuchenne.

%%K t06-c4-08-1 p
8. --- \VUgras{Miotła} \textsuperscript{(34)} i \VUgras{miotełka z piór}
\textsuperscript{(35)} są w kącie. Kucharka dopiero co używała \VUgras{kloca
drewnianego} \textsuperscript{(36)}, zostawiła go obok okna.

%%K t06-c4-09-1 p
9. --- \VUgras{Ręcznik} \textsuperscript{(37)} wisi na ścianie, obok \VUgras{lejka}
\textsuperscript{(38)}. W tej części kuchni odłożone są \VUgras{ruszt}
\textsuperscript{(39)}, \VUgras{tasak} \textsuperscript{(40)} i \VUgras{deska do
siekania} \textsuperscript{(41)}. Dalej, nad tasakiem, na \VUgras{półce}
\textsuperscript{(42)}, są \VUgras{garnki} \textsuperscript{(43)}, \VUgras{garnek do
gotowania} \textsuperscript{(44)} ze swoją \VUgras{pokrywką} \textsuperscript{(45)} i
\VUgras{kociołek} \textsuperscript{(46)}. \VUgras{Patelnia} \textsuperscript{(47)} wisi
obok.

%%K t06-c4-10-1 p
10. --- Tylko miedziany \VUgras{kran} \textsuperscript{(48)} rzuca blask w tym kącie.
\VUgras{Zlew} \textsuperscript{(49)}, na którym stoi gruby \VUgras{dzban}
\textsuperscript{(50)}, jest zapewne bardzo wygodny ze swoją małą szafką, w której
przechowuje się \VUgras{beczułkę octu} \textsuperscript{(51)} zaopatrzoną w swój
\VUgras{kranik} \textsuperscript{(52)}, \VUgras{pudło na odpadki} \textsuperscript{(53)}
i \VUgras{brudne talerze} \textsuperscript{(54)}.

%%K t06-tit-9 sub
\VUsaut{4.0mm}
\VUpk{10.2pt}{40}{Inne przedmioty umieszczone w kuchni.}
\VUsaut{3.0mm}

%%K t06-c4-11-1 p
11. --- \VUgras{Balia} \textsuperscript{(55)}, w której myje się \VUgras{warzywa}
\textsuperscript{(58)}, i żeliwny \VUgras{kocioł} \textsuperscript{(56)} postawiony na
\VUgras{posadzce} \textsuperscript{(57)} mają zapewne także swoje miejsce w tej szafce.

%%K t06-c4-12-1 p
12. --- Kucharce z pewnością nie brakuje pieców, bo oto jeszcze, na rogu stołu,
\VUgras{kuchenka gazowa} \textsuperscript{(59)}, na której grzeje się woda w małym
rondlu. \VUgras{Para} \textsuperscript{(60)}, która z niego ucieka, wskazuje nam, że
zapewne się gotuje. Prawdopodobnie ta woda posłuży do kawy, bo oto na drugim końcu
stołu \VUgras{dzbanek do kawy} \textsuperscript{(64)} i \VUgras{młynek do kawy}
\textsuperscript{(63)}.

%%K t06-c4-13-1 p
13. --- Kuchenka jest połączona z \VUgras{palnikiem gazowym} \textsuperscript{(62)}
długim \VUgras{wężykiem gumowym} \textsuperscript{(61)}.

%%K t06-c4-14-1 p
14. --- \VUgras{Posługaczka} \textsuperscript{(66)}, która przychodzi pomagać kucharce,
przygotowuje warzywa na obiad. Kiedy zostaną oczyszczone i umyte, położy je w
\VUgras{misce} \textsuperscript{(65)}, czekając na godzinę ich gotowania.

%%K t06-tit-10 sub
\VUsaut{4.0mm}
{\centering\fontsize{10.2pt}{10.2pt}\selectfont\textls[40]{\scshape Łazienka.}
 \textit{(Pokój kąpielowy.)}\par}
\VUsaut{3.0mm}

%%K t06-c5-01-1 p
1. --- Pozostaje nam do zrobienia tylko jedna niedyskrecja, żeby poznać główne pokoje
mieszkania: zobaczyć urządzenie łazienki. To nie zabierze dużo czasu, bo nie jest
wielka, i szybko dowiemy się, że \VUgras{wanna} \textsuperscript{(1)} ma dobry
\VUgras{piecyk kąpielowy} \textsuperscript{(2)}, że \VUgras{mydelniczka}
\textsuperscript{(3)} jest w zasięgu ręki kąpiącego się i że \VUgras{wieszak na
ręczniki} \textsuperscript{(5)} z pewnością pozwoli łatwo wysuszyć \VUgras{ręczniki}
\textsuperscript{(4)}.

%%K t06-c5-02-1 p
2. --- Ten pokój ma ściany pokryte \VUgras{kafelkami} \textsuperscript{(6)}, które
pozwoliły zainstalować \VUgras{natrysk} \textsuperscript{(7)} bez żadnej obawy, że
woda, która rozpryskuje się przy jego używaniu, uszkodzi mury. Cynkowa \VUgras{miednica}
\textsuperscript{(8)}, która służy do kąpieli nóg, uzupełnia umeblowanie.
\VUsaut{6.0mm}
\VUfilet{22mm}
